% SPDX-License-Identifier: CC-BY-4.0
% Copyright (c) 2026 Martin Schröder <info@swedishembedded.com>

\section{Reproducibility}\label{sec:repro}

\paragraph{Verification.} \code{make -C papers/splinter-brain-sven-flywheel verify} first runs the
specification of the pin rule, then re-runs every structure check against the commits in
\code{data/commits.json}, then \code{scripts/build_data.py --check}, which resolves each measured,
computed, registered and planned claim to its record, requires the claim text to contain the value,
and requires every ledger row to be cited in the text and every citation to be in the ledger.
The ledger holds \mLedgerClaims{} claims, \mLedgerStructure{} of them structure claims backed by
\mLedgerChecks{} checks. A pin must be contained in a remote-tracking ref of its repository, so a
fresh clone of the three remotes reproduces the result; the verification was run on such a clone
of the remotes, with the sibling repositories given by \code{--sven} and \code{--brain}. The
splinter pin is the reachable twin of an earlier pin that no branch contained.
\code{make check} includes this verification as \code{check/papers}. The bibliography entries
were checked against the archive's abstract pages on 2026-10-16.

\paragraph{Experiments.} The exact commands are in \code{scripts/run_experiments.sh}; raw records
are copied into \code{data/runs} by \code{scripts/archive_run.sh} with this machine's paths replaced.
The first wave ran a Splinter binary built from the main checkout at the time (Splinter
\code{ce42a55} plus uncommitted repairs, Sven \code{02adbda7}, Brain \code{dc50b7cc}; binary sha256
prefix \code{f91620c5}); its records are in \code{data/runs/e1} and \code{data/runs/e2}. The second wave
ran binaries of Splinter's main branch at the commit before this paper's change plus the repairs of
\cref{sec:defects} in the order listed there (the E1 reruns and the registered E2 run the first three
repairs and the serve-bank repair; the repaired-verifier rerun and E2d also the verifier repair), against Sven
at its local main \code{ebd169a2} and a Brain release binary built from \code{f9f909f7}
(\code{--features brain/study}), with a CUDA 13 toolkit on the card's host. Splinter's models run in its process
(Vulkan adapter line in the logs); \code{brain serve} in the gate's serve check ran with the CUDA backend.
Models: Qwen3-0.6B or Qwen3-1.7B (policy), Qwen3-4B (generator, teacher; Qwen3-8B in E1c), on a shared NVIDIA GH200, so wall times are not
comparable between runs. The seeds drawn by the runs are in the raw records (denoise seed in
\code{tasks.json}, fact seed in \code{build.json}, training seeds in each \code{learn-g*.stderr}). The protocol was
fixed before the runs and extended before each later one (\code{prereg/protocol.md}); the pilots that chose the
training flags are described there and are not part of any claim. Deviations: E1 ran without a time limit; the registered E2 success
criterion was not met; a repaired run, a post hoc run and a post hoc E1 variant are reported beside it and labelled;
the E2d run was assembled from three invocations of the same script
(\code{run} until its first served probe failed on an authentication bug of the script, repaired in \code{finish}, whose
served probes were stopped under load; the rollback of its first release; and \code{second\_generation.py}, the second \code{learn} on a copy
of its state taken before the rollback), all on the state
the first produced, nothing retried.

\paragraph{Training qualification and throughput.} The logs in \code{data/runs/qual} are Brain's
qualification logs at commit \mEFourCommit{}; each log's header records its command. They contain
only the runs that finished; runs that were planned and not run are not reported as results. The
throughput probe in \code{data/runs/throughput} is a copy of a log kept with the correction plan;
its header says what the roadmap records about its command and revision, and that the log itself
carries no commit.

\paragraph{Coding-task experiments.} None has been run. The registration is
\code{prereg/coding-tasks.md}; results will be rows of \code{data/results.csv} with the commands,
commits, harness profile and seeds of the run, and the tables of \cref{sec:pending} are generated
from them.

\paragraph{Build.} \code{make papers/splinter-brain-sven-flywheel/pdf} builds the paper to
\code{target/papers}; the tables and numbers are generated from \code{data/runs} into the
build directory.
